%% kapa: exam=1 round=33 year=2022 period=2 form=A subject=accounting subject_ko=회계학 seq=24 number=64
%% source: Q-Net 감정평가사 기출문제 (cst00309) artlSeq=5209997 file=raw/1008/5209997/2022년도 제33회 감정평가사 1차시험 2교시(감정평가 관계법규, 회계학).pdf
%% conversion: kordoc 4.24.51 -> tools/gen1.py (auto); review=proofread-sonnet; flags=none
\begin{question}{64}
퇴직급여제도에 관한 설명으로 옳지 않은 것은?
\begin{choices}
  \item 확정기여제도에서는 종업원이 보험수리적위험(급여가 예상에 미치지 못할 위험)과 투자위험(투자자산이 예상급여액을 지급하는데 충분하지 못할 위험)을 실질적으로 부담한다.
  \item 확정기여제도에서는 기여금의 전부나 일부의 납입기일이 종업원이 관련 근무용역을 제공하는 연차보고기간 말 후 12개월이 되기 전에 모두 결제될 것으로 예상되지 않는 경우를 제외하고는 할인되지 않은 금액으로 채무를 측정한다.
  \item 확정급여채무의 현재가치와 당기근무원가를 결정하기 위해서는 예측단위적립방식을 사용하며, 적용할 수 있다면 과거근무원가를 결정할 때에도 동일한 방식을 사용한다.
  \item 확정급여제도에서 기업이 보험수리적위험(실제급여액이 예상급여액을 초과할 위험)과 투자위험을 실질적으로 부담하며, 보험수리적 실적이나 투자실적이 예상보다 저조하다면 기업의 의무가 늘어날 수 있다.
  \item 퇴직급여채무를 할인하기 위해 사용하는 할인율은 보고기간 말 현재 그 통화로 표시된 국공채의 시장수익률을 참조하여 결정하고, 국공채의 시장수익률이 없는 경우에는 보고기간 말 현재 우량회사채의 시장수익률을 사용한다.
\end{choices}
\end{question}
\answer{5}
